\documentclass[10pt,a4paper]{article}
\usepackage[brazilian]{babel}
\usepackage{../../styles/resume-style}

\begin{document}

\begin{center}
{\Large\bfseries Yago Henrique Teixeira}\\[3pt]
\textbf{Desenvolvedor iOS Sênior | Swift | SwiftUI | Arquitetura Mobile}\\[3pt]
Belo Horizonte, MG \quad | \quad (31) 99834-8635 \quad | \quad
\href{mailto:yago.ht@icloud.com}{yago.ht@icloud.com} \quad | \quad
\href{https://www.linkedin.com/in/yagohtx/}{linkedin.com/in/yagohtx}
\end{center}

\section{Resumo Profissional}

Desenvolvedor iOS com mais de 6 anos de experiência em desenvolvimento de software, especializado em Swift, SwiftUI e arquitetura mobile. Tenho experiência na modernização de aplicações legadas, criação de componentes reutilizáveis, modularização de código e automação de fluxos de engenharia com IA. Na AB InBev, atuo em um aplicativo B2B de e-commerce utilizado em múltiplos países, contribuindo para decisões de arquitetura e evolução de módulos compartilhados.

\section{Competências Técnicas}

\textbf{Linguagens:} Swift, Java, JavaScript, SQL

\textbf{iOS e plataformas Apple:} SwiftUI, UIKit, Combine, async/await, Core Data, SwiftData, Realm, Deep Links

\textbf{Arquitetura:} MVVM, The Composable Architecture (TCA), Clean Architecture, Coordinator, injeção de dependências (Swinject), arquitetura modular, ViewCode

\textbf{Testes:} XCTest, Quick, Nimble, SnapshotTesting, ViewInspector, testes unitários e de snapshot

\textbf{Observabilidade e analytics:} Segment, Firebase Remote Config, Crashlytics, New Relic

\textbf{Ferramentas e CI/CD:} Git, Swift Package Manager (SPM), CocoaPods, Azure Pipelines, Fastlane, Mockoon

\textbf{Engenharia assistida por IA e automação:} Cursor, agentes de IA, Skills, Model Context Protocol (MCP), Azure DevOps, Jira, automação de workflows

\textbf{Backend e dados:} Spring Boot, Node.js, APIs REST, PostgreSQL, MongoDB, Apache Kafka, WebSockets

\section{Experiência Profissional}

\job{Desenvolvedor iOS}{Jun 2022 -- Nov 2026}{AB InBev}{Belo Horizonte, MG}

\begin{itemize}

\item Desenvolvo e evoluo um aplicativo iOS B2B de e-commerce utilizado em múltiplos países, conectando varejistas ao catálogo de produtos da AB InBev.

\item Implementei componentes reutilizáveis em SwiftUI, filtros e novas experiências de interface integradas a um Back-end for Front-end (BFF), incluindo o mapeamento de respostas para a camada de apresentação e o tratamento de erros com async/await.

\item Liderei a migração de funcionalidades de descoberta de produtos para SwiftUI e The Composable Architecture (TCA), isolando estado e responsabilidades e facilitando a evolução da aplicação.

\item Desenvolvi busca offline com Realm e, posteriormente, SwiftData, incluindo buscas recentes, deep links e navegação orientada a estado.

\item Instrumentei eventos no Segment ao longo de jornadas importantes, como navegação, filtros, buscas e compras, ampliando a observabilidade do produto.

\item Desenvolvo testes unitários e de snapshot para reduzir regressões, além de utilizar New Relic e Crashlytics na investigação de problemas de performance e falhas em produção.

\item Entrego funcionalidades para diferentes mercados utilizando localização, Firebase Remote Config, AppLanga, configurações específicas por país e testes A/B com Optimizely.

\item Atuo em um monorepo modular com nove módulos compartilhados entre diferentes equipes. Também implementei o Mockoon em um aplicativo interno de demonstração, reduzindo a dependência do backend durante testes locais.

\item Estruturei fluxos de engenharia assistidos por IA com Cursor, Skills e MCP integrados ao Azure DevOps e Jira, apoiando refinamentos, implementação, documentação e criação de pull requests.

\item Atuo como referência técnica em SwiftUI e arquitetura mobile, contribuindo em code reviews, discussões técnicas e decisões de arquitetura entre equipes.

\end{itemize}

\job{Desenvolvedor iOS Freelancer}{Jan 2024 -- Dez 2024}{Rhitmo Tech}{Belo Horizonte, MG}

\begin{itemize}

\item Liderei a modernização completa de um aplicativo iOS de controle de ponto, migrando uma base monolítica em UIKit, sem arquitetura definida, para MVVM com SwiftUI.

\item Migrei todas as telas do aplicativo para SwiftUI, preservando as regras de negócio existentes e reduzindo significativamente a complexidade de manutenção da aplicação.

\item Estruturei um sistema de componentes reutilizáveis inspirado em Atomic Design, padronizando interfaces e reduzindo duplicação de código.

\item Reestruturei a camada de networking, substituindo uma implementação centralizada composta por métodos extensos por uma arquitetura modular com responsabilidades bem definidas.

\item Reestruturei persistência, navegação e regras de negócio com Core Data, além de implementar testes unitários para as funcionalidades principais.

\item Reconstruí regras de negócio por meio de engenharia reversa das APIs e documentação incompleta, transformando um sistema legado sem arquitetura em uma base sustentável.

\item Atuei como referência técnica do projeto, mentorando dois desenvolvedores, participando da contratação técnica da equipe e conduzindo decisões de arquitetura.

\end{itemize}

\job{Desenvolvedor Java}{Nov 2020 -- Jun 2022}{Avenue Code}{Belo Horizonte, MG}

\begin{itemize}

\item Desenvolvi aplicações web utilizando Java e Spring Boot, implementando APIs REST, regras de negócio e integrações com SQL e MongoDB.

\item Trabalhei em fluxos assíncronos utilizando Apache Kafka e desenvolvi suítes de testes automatizados com Cucumber (BDD).

\end{itemize}

\job{Estagiário de Desenvolvimento Java}{Mai 2020 -- Nov 2020}{Avenue Code}{Belo Horizonte, MG}

\begin{itemize}

\item Desenvolvi serviços backend em Java, Spring Boot e Node.js para aplicações internas.

\item Implementei comunicação assíncrona com Apache Kafka para distribuição de eventos em uma plataforma social corporativa.

\end{itemize}

\section{Formação Acadêmica}

\textbf{Bacharelado em Ciência da Computação} \hfill 2017 -- 2020

Centro Universitário de Belo Horizonte (UniBH) \hfill CR: 91,29/100

\section{Informações Adicionais}

\textbf{Idiomas:} Português nativo; inglês fluente, utilizado diariamente em reuniões, documentação e colaboração com equipes internacionais.

\textbf{Práticas de trabalho:} Scrum, code review, pair programming, refinamentos técnicos e Planning Poker.

\end{document}

